

Proviamo a dimostrare \(\vdash (\exists x, (P(x)\lor R(x)) )\to (\exists x, P(x) \lor \exists x, R(x))\).

Qui ancora per \(\to I\) dobbiamo dimostrare \(\exists x, P(x) \lor \exists x, R(x)\) potendo assumere \(\exists x, (P(x)\lor R(x))\). Quando andrò ad usare l'assuzione \(\exists x, (P(x)\lor R(x))\) avrò a disposizione un ulteriore assunzoine scaricata che è \(P(y) \lor R(y)\)

Ricordiamoci di dover dimostrare una formula con \(\lor\) come operatore principale, quindi potremmo pensare di ottenerla tramite \(\lor I\), ad esempio da \(\exists x\st P(x)\), che è ottenibile da \(P(y)\) con \(\exists I\). Purtroppo però nessuna delle due assunzioni che posso scaricare mi permettono di scaricare \(P(y)\). Ma ho a disposizione \(P(y) \lor R(y)\), e di fatto quello che ho fatto è assumere un disgiunto di tale \(\lor\). Se quindi posso dimostrare \(\exists x\st P(x) \lor \exists x\st R(x)\) anche dall'altro disgiunto, allora con \(\lor I\) potrò scaricare entrambe le assunzioni. In realtà posso con lo stesso ragionamento. A questo punto è fatta perché abbiamo ottenuto \(\exists x, P(x) \lor \exists x, R(x)\) dalla sola assunzione \((P(x)\lor R(x))\). Abbiamo però già detto che questa assunzione la possiamo scaricare mediante \(\exists E\) dalla nostra assunzione iniziale, se tutti i vincoli sono rispettati. In questo caso \(\phi = P(x) \lor R(x)\), in cui non compare \(y\). Inoltre non compare libera in \(\exists x, P(x) \lor \exists x, R(x)\), e non compare in nessuna assunzione non scaricata. Quindi tutti i vincoli sono rispettati, e quindi possiamo scaricare l'assunzione \(P(y) \lor R(y)\) per ottenere \(\exists x\st P(x) \lor \exists x\st R(x)\). A questo punto con \(\to I\) otteniamo la tesi.

Passiamo al verso opposto, ovvero \((\exists x, P(x) \lor \exists x, R(x))\to (\exists x, (P(x)\lor R(x)) )\).

Chiaramente come al solito l'operatore principale è \(\to\), quindi iniziamo con \(\to I\) dovendo dimostrare \(\exists x, (P(x)\lor R(x))\) potendo assumere \(\exists x, P(x) \lor \exists x, R(x)\). Tale assunzione è una disgiunzione, quindi potrà regalarmi come assunzioni che scaricherà i due disgiunti, ovvero \(\exists x\st P(x)\) e \(\exists x\st R(x)\). Ancora una volta, questi esistenziali per \(\exists I\) ci regalano \(P(y)\) e \(R(y)\) come assunzioni scaricabili.

A questo punto in parallelo possiamo da \(P(y)\) e \(R(y)\) ottenere \(P(y)\lor R(y)\) con \(\lor I\), e da questi con \(\exists I\) ottenere \(\exists x, (P(x)\lor R(x))\). Chiaramente non avrei potuto usare direttamente \(\exists E\) per scaricare subito \(P(y)\) perché \(y\) è libero in \(P(y)\). Applicando però prima \(\exists I\) ottengo \(\exists x\st P(x)\) rimuovendo le occorrenze di \(y\) se non in \(P(y)\), e quindi ora posso scaricare \(P(y)\) con \(\exists E\) ottenendo ancora \(\exists x, (P(x)\lor R(x))\). Avendo ottenuto \(\exists x, (P(x)\lor R(x))\) da entrambi i disgiunti possiamo con \(\lor I\) ottenere \(\exists x, (P(x)\lor R(x))\) da \(\exists x, P(x) \lor \exists x, R(x)\). A questo punto con \(\to I\) otteniamo la tesi.

Proviamo ora a dimostrare \( \forall x, (A(y) \to B(x)) \vdash A(y) \to \forall x, B(x)\).

Per dimostrare \(A(y) \to \forall x, B(x)\) dovremo usare \(\to I\), e quindi dovremo dimostrare \(\forall x, B(x)\) potendo assumere \(A(y)\). A questo punto per dimostrare \(\forall x, B(x)\) dovremo usare \(\forall I\), e quindi dovremo dimostrare \(B(z)\) potendo assumere \(A(y)\). Per non avere problemi con i vincoli di applicabilità di \(\forall I\) è utile usare variabili sempre nuove quando si può in modo da rendere più facile verificare i vincoli. Ora \(B(z)\) la potrei ottenere per contraddizione, ma con \(\neg (B(z))\) ce ne potremo fare ben poco, quindi conviene usare \(\forall E\) per ottenere \(A(y) \to B(z)\), che posso usare perché \(z\) è libero per \(x\) in \(A(y) \to B(x)\). A questo punto con \(A(y)\) ottengo \(B(z)\) che ci serve per concludere la dimostrazione. 

Passiamo all'altro verso, ovvero \(A(y) \to \forall x, B(x) \vdash \forall x, (A(y) \to B(x))\).

Volendo ottenere una formula quantificata, ragionevolmente dovremo usare \(\forall I\). Per farlo dobbiamo dimostrare \(A(y) \to B(z)\) potendo assumere \(A(y) \to \forall x, B(x)\). Tale \(z\) non dovrà apparire in nessuna assunzione non scaricata.
Ora \(A(y) \to B(z)\) è una formula con \(\to\) come operatore principale, quindi per ottenerla dovremo usare \(\to I\). Per farlo dobbiamo dimostrare \(B(z)\) potendo assumere \(A(y)\). Ma adesso con questa assunzione e l'assunzione \(A(y) \to \forall x, B(x)\) possiamo ottenere \(\forall x, B(x)\) con \(\to E\), e da questa con \(\forall E\) otteniamo \(B(z)\) che ci serve per concludere la dimostrazione.

Passiamo ora a un esempio più complesso, ovvero \(\forall x, P(x) \dashv \vdash \neg \exists x\st \neg P(x)\). Essendo che questa è la generalizzazione di De Morgan al prim'ordine, mi aspetto che la struttura di dimostrazione sia simile a quella di De Morgan.

Iniziamo con la direzione \(\forall x, P(x) \dashv \neg \exists x\st \neg P(x)\).

Qui come unica assunzione abbiamo \(\neg \exists x\st \neg P(x)\), che non è molto usabile se non avendo la sua versione positiva.

Notiamo però che \(\forall x, P(x)\) è ottenibile tramite \(\forall I\) da \(P(y)\). Ora l'unico modo ragionevole per ottenere \(P(y)\) è tramite \(\RAA\). Ora \(\neg P(y)\) è un assunzione che potremo scaricare, e ci serve dimostrare una contraddizione. Ora da \(\neg P(y)\) è ottenibile \(\exists x\st \neg P(x)\) con \(\exists I\), e da questa con \(\neg E\) otteniamo finalmente \(\bot\) che ci serve per concludere la dimostrazione.

Possiamo vedere l'analogo proposizoinale precedentemente citato. Consideriamo ad esempio \(\neg (\neg A \lor \neg B)\vdash A \land B\). La dimostrazione qui è per \(\land I\) da \(A\) e \(B\). Per ottenere \(A\) dobbiamo usare \(\RAA\), e quindi assumere \(\neg A\) per ottenere una contraddizione. Ora da \(\neg A\) è ottenibile \(\neg A \lor \neg B\) con \(\lor I\), e da questa con \(\neg E\) otteniamo \(\bot\) che ci serve per concludere la dimostrazione. La dimostrazione di \(B\) è del tutto analoga.

Possiamo notare che la dimostrazione è del tutta analoga a quella appena vista con i quantificatori, in cui si scambiano le regole per \(\forall\) con quelle per \(\land\), e le regole per \(\exists\) con quelle per \(\lor\).

Passiamo dunque all'altro verso, ovvero \(\neg \exists x\st \neg P(x) \vdash \forall x, P(x)\).

La conclusione che dobbiamo dimostrare è una formual negata, quindi è ragionevole usare \(\neg I\), che ci chiede di dimostrare \(\bot\) potendo assumere \(\exists x\st \neg P(x)\). A questo punto da \(\forall x, P(x)\) è ottenibile \(P(y)\) con \(\forall E\), che ci da \(\bot\) con \(\neg E\) assumendo \(\neg P(y)\). Per scaricare questa assunzione possiamo usare \(\exists E\) con \(\phi = \neg P(x)\), che non contiene \(y\), e non contiene \(y\) libera in nessuna assunzione non scaricata, e non contiene \(y\) libera in \(\bot\). Quindi tutti i vincoli sono rispettati, e quindi possiamo scaricare l'assunzione \(\neg P(y)\) per ottenere \(\bot\) che ci serve per concludere la dimostrazione.

Possiamo ora mostrare che \(\neg \exists x \st P(x)\vdash \exists x \st \neg P(x)\).

Per dimostrare \(\exists x \st \neg P(x)\) è ragionevole usare \(\exists I\), che ci chiede di dimostrare \(\neg P(y)\). Ora essendo una negazione, e non avendola nelle assunzioni, è ragionevole usare \(\neg I\), che ci chiede di dimostrare \(\bot\) potendo assumere \(P(y)\). Ora da \(P(y)\) è ottenibile \(\exists x \st P(x)\) con \(\exists I\), e da questa con \(\neg E\) otteniamo \(\bot\) che ci serve per concludere la dimostrazione.

Proviamo a vedere se l'altro verso è dimostrabile. Vogliamo dimostrare \(\neg \exists x \st P(x)\) e l'unica assunzione che abbiamo è \(\exists x \st \neg P(x)\). Ora per dimostrare \(\neg \exists x \st P(x)\) è ragionevole usare \(\neg I\), che ci chiede di dimostrare \(\bot\) potendo assumere \(\exists x \st P(x)\). Ora da queste due assunzioni si potrebbe pensare di scaricare \(P(y)\) e \(\neg P(y)\) per permetterci di dedurre \(\bot\). Purtroppo però \(y\) quando si scarica \(\neg P(y)\) tramite \(\exists E\) è libera in \(P(y)\), che non è scaricata. Si è quindi violato un vincolo di applicabilità, e quindi non è possibile concludere la dimostrazione. Questo esempio è particolarmente sottile, poiché \(P(y)\) verrà scaricata in seguito, e quindi si potrebbe pensare che non ci siano problemi perché tutte le assunzioni con \(y\) sono scaricate. Purtroppo però non sono scaricate tutte assieme, nel momento in cui si scarica \(\neg P(y)\) abbiamo \(P(y)\) ancora non scaricata, e quindi non è possibile concludere la dimostrazione.

Infine, vediamo la dimostrazoine di \(\vdash \exists x \st (C(x)\to \forall y, C(y))\).

Per dimostrare \(\exists x \st (C(x)\to \forall y, C(y))\) è ragionevole usare \(\exists I\), che ci chiede di dimostrare \(C(z) \to \forall y, C(y)\). Ora \(C(z) \to \forall y, C(y)\) è una formula con \(\to\) come operatore principale, quindi per ottenerla è ragionevole usare \(\to I\), che ci chiede di dimostrare \(\forall y, C(y)\) potendo assumere \(C(z)\). 

Ora \(\forall y, C(y)\) è ottenibile sia per contraddizione che per \(\forall I\). Iniziamo a provare la contraddizione. Per farlo potremmo assumere \(\neg C(z)\) da scaricare. Per scaricarla possiamo utilizzare \(\RAA\), ottenendo una contraddizione tra \(\exists x \st (C(x)\to \forall y, C(y))\) e \(\neg \exists x \st (C(x)\to \forall y, C(y))\). Ora abbiamo ottenuto nuovamente \(C(z)\), da cui possiamo ottenere \(\forall y, C(y)\) con \(\forall I\), che con \(\to I\) ci ridà \((C(x)\to \forall y, C(y))\) che con \(\exists I\) ci da la tesi. Ci manca però da scaricare \(\neg \exists x \st (C(x)\to \forall y, C(y))\) che possiamo farlo con un ultima applicazione di \(\RAA\) che ci da la tesi.

\chapter{Lezione 14/05}

Vediamo altri schemi di deduzione in \(\FOL\) con deduzione naturale. 

Nella logica proposizoinale se deduciamo dalle stesse assunzioni e da \(P\) un certo \(\psi\) e dalle stesse assunzioni e \(\neg P\) otteniamo \(\psi\), allora possiamo dedurre \(\psi\) scaricando le assunzioni di \(P\) o \(\neg P\) tramite \(\lor E\) con \(P \lor \neg P\) che è una tautologia quindi non richiede assunzioni. Quest'ultima si può dimostrare tramite \(\neg E, \lor i, \neg i \neg e, \lor i\) dal basso verso l'altro.

Un altro modo sempre avendo le due deduzioni di \(\psi\) da \(P\) e da \(\neg P\) è quello di usare \(\neg E\) alla deduzione da \(P\) assumendo \(\neg \psi\), per poi scaricare \(P\) tramite \(\neg I\), dal \(\neg P\) ottenuto potrò usare la seconda deduzione ottenendo ancora \(\psi\) che ancora con \(\neg E\) posso usare per scaricare \(\neg \psi\) tramite \(\RAA\) e quindi avere \(\psi\).

Un esempio dell'applicazioni di questi schemi la possiamo trovare in \(p\to q \vdash \neg p \lor q\).
AGGIUNGI DEDUZIONI

Nel prim'ordine è più sottile, perché se una dimostrazione ha la prima struttura sarà sempre possibile ottenere una dimostrazione con la seconda struttura, ma non è detto che sia possibile il contrario, poiché scomponendo la dimostrazione in due rami si mantengano i vincoli di applicabilità delle regole. Questo si vede bene nell'ultima dimostrazione vista.

TODO Ragiona sul fatto di chi è \(\pi_1\) e chi \(\pi_2\) e che nel primo schema \(C(z)\) non è assunzione, nel secondo si, e questo viola vincolo \(z\) non compare libero in nessuna assunzione non scaricata.

Andiamo ora a introdurre \(=\) come simbolo logico, e con quindi delle vere e proprie regole di deduzione naturale per \(=\).
Come già detto l'uguaglianza come simbolo logico è solo un opzione, perché si può anche fare tramite assiomi, ovvero formule date per valide. In particolare si sta arricchendo il linguaggio con un concetto di Teoria, ovvero un insieme di formule valide in ogni interpretazione, che è detto rappresentabile se esiste un insieme di formule che deducono tutte le formule della teoria, detto appunto insiemi di assiomi. Tratteremo questo concetto in maniera più approfondita in seguito, ma possiamo pensare subito a un esempio di teoria semplice, ovvero la teoria dell'uguaglianza. In particolare dovrà avere
\begin{itemize}
  \item \(\forall x, x=x\)
  \item \(\forall x_1,\ldots, \forall x_n, \forall y_1, \ldots, \forall y_n, (\bigwedge_{i=1}^n x_i = y_i) \to (f(x_1, \ldots, x_n) = f(y_1, \ldots, y_n))\) per ogni \(f \in \FuncSyms\) con \(\arty(f) = n\).
  \item \(\forall x_1,\ldots, \forall x_n, \forall y_1, \ldots, \forall y_n, (\bigwedge_{i=1}^n x_i = y_i) \to (P(x_1, \ldots, x_n) \to P(y_1, \ldots, y_n))\) per ogni \(P \in \PredSyms\) con \(\arty(P) = n\).
\end{itemize}

Per lavorare in questa teoria andrebbe aggiunta alla definizione di deduzione naturale la possibilità di lasciare come assunzione non scaricata una qualsiasi di queste formule.

Ci si aspetterebbe che il terzo assioma abbia una doppia implicazione, ma non viene messo perché deducibile dall'implicazione semplice. Se volessimo infatti dimostrare che \(t=s \to (P(t)\leftrightarrow P(s))\) da questo terzo assioma. Infatti da \(\forall x, \forall y, x=y \to (P(x) \to P(y))\) possiamo applicare due volte \(\forall E\) istanziando prima \(x\) con \(t\) e \(y\) con \(s\) per ottenere \(t=s \to (P(t) \to P(s))\). Per l'altro verso sarà sufficiente istanziare prima \(x\) con \(s\) e \(y\) con \(t\) per ottenere \(s=t \to (P(s) \to P(t))\). Ora mancherebbe la simmetria di \(=\) per concludere la dimostrazione, ma essendo \(= \in \PredSyms\), il terzo assioma insieme alla riflessività, come vedremo più avanti, ci permetterà di dimostrare la simmetria di \(=\), e quindi con \(\to I\) otteniamo \(t=s \to (P(t) \leftrightarrow P(s))\).


Se volessimo invece usare l'approccio con regole di deduzione naturale, possiamo introdurre una regola per ogni assioma, ovvero
\[\infer[Refl]{t=t}{}\]
\[\infer[Sub_f]{\subs{f(x_1, \ldots, x_n)}{x_1,\ldots,x_n}{t_1,\ldots, t_n} = \subs{f(x_1, \ldots, x_n)}{x_1,\ldots,x_n}{s_1,\ldots,s_n}}{t_1 = s_1 & \ldots & t_n = s_n}\]
\[\infer[Sub_P]{\subs{P(x_1, \ldots, x_n)}{x_1,\ldots,x_n}{s_1,\ldots,s_n}\to\subs{P(x_1, \ldots, x_n)}{x_1,\ldots,x_n}{t_1,\ldots, t_n}}{t_1 = s_1 & \ldots & t_n = s_n}\]

Proviamo ora a dimostrare \(x=y \vdash y=x\). Usando \(Refl\) otteniamo \(x = x\). A questo pundo definendo \(P(x_1,x_2) \eqdef x_1 = x_2\) possiamo applicare \(Sub_P\) ottenendo \(x=x \to y = x\) e a questo punto con \(x=x\) ottenuto da \(Refl\) otteniamo \(y=x\) che ci serve per concludere la dimostrazione.

Possiamo anche mostrare la transitività, ovvero \(x=y,y=z\vdash x =z\). Chiaramente useremo la stessa \(P\) definita prima. Dalla simmetria di prima sappiamo che \(x=y \vdash y=x\) e a questo punto con \(y=z\) e \(Sub_P\) otteniamo \(y=y\to x=z\), da cui per \(\to E\) possiamo ottenere \(x=z\), assumendo \(y=y\) che è ottenibile da \(Refl\).

Le regole di sostituzione possono esseren generalizzate. \(Sub_f\) può essere infatti generalizzata a qualsiasi termine \(t\) con:
\[\infer[Sub_t]{\subs{t}{x_1,\ldots,x_n}{s_1,\ldots,s_n} = \subs{t}{x_1,\ldots,x_n}{t_1,\ldots,t_n}}{s_1=t_1 & \ldots & s_n=t_n}\]

Chiaramente questa nuova regola è derivabile dalla prima, poiché nonostante la seconda ammetta anche termini innestati è possibile applicare la regola sui termini atomici che compaiono in \(t\) singolarmente più volte in modo da ricostruire la regola generalizzata. Si può chiaramente dimostrare per induzione strutturale sulla struttura dei termini.

Sulla falsa riga di questa potremmo generalizzare anche \(Sub_P\) a qualsiasi formula \(\phi\) con:
\[\infer[Sub_\phi]{\subs{\phi}{x_1,\ldots,x_n}{s_1,\ldots,s_n}\to\subs{\phi}{x_1,\ldots,x_n}{t_1,\ldots,t_n}}{s_1=t_1 & \ldots & s_n=t_n}\]

Chiaramente anche in questo caso, per induzione strutturale sulle formule, possiamo dimostrare che questa regola è derivabile da \(Sub_P\). Il caso base è chiaramente \(Sub_P\). Tra i casi induttivi ce ne sono alcuni interessanti. Consideriamo ad esempio il caso in cui \(\phi = \neg \psi\). Ovviamente a \(\psi\) si applicano ipotesi induttiva. A questo punto invertendo tutte le uguaglianze per simmetria per ipotesi induttiva potremmo ottenere \(\subs{\psi}{x_1,\ldots,x_n}{s_1,\ldots,s_n} \to \subs{\psi}{x_1,\ldots,x_n}{t_1,\ldots,t_n}\). Ora assumendo \(\subs{\psi}{x_1,\ldots,x_n}{s_1,\ldots,s_n}\) otteniamo \(\subs{\psi}{x_1,\ldots,x_n}{t_1,\ldots,t_n}\) con \(\to E\), da cui assumendo \(\neg (\subs{\psi}{x_1,\ldots,x_n}{t_1,\ldots,t_n}) = \subs{\neg\psi}{x_1,\ldots,x_n}{t_1,\ldots,t_n}\) otteniamo \(\bot\) con \(\neg E\). A questo punto con \(\neg I\) otteniamo \(\subs{\neg\psi}{x_1,\ldots,x_n}{s_1,\ldots,s_n} \) e infine con \(\to I\) otteniamo \(\subs{\neg\psi}{x_1,\ldots,x_n}{s_1,\ldots,s_n} \to \subs{\neg\psi}{x_1,\ldots,x_n}{t_1,\ldots,t_n}\) che è proprio quello che ci serve per concludere il caso induttivo.

A questo punto andiamo a dimostrare che le condizioni di applicabilità delle regole di quantificazione sono sufficienti, oltre che necessarie. In particolare, per fare ciò vedremo che il calcolo di deduzione naturale così definito è corretto, \(\Gamma \vdash \phi \implies \Gamma \models \phi\). Come fatto nella logica proposizionale, procederemo per induzione strutturale sull'albero di deduzione. Il caso base è chiaramente l'albero di altezza \(0\), che è il caso di \(\phi \in \Gamma\). Chiarmente in questo caso \(\Gamma \models \phi\) per definizione di soddisfacibilità. Passando al caso induttivo, supponiamo di avere una deduzione con radice \(\phi\) e ipotesi \(\Gamma\) e che tutti i sottoalberi della radice sono tali che se hanno ipotesi \(\Gamma'\) e radice \(\psi\) allora \(\Gamma' \models \psi\). Chiaramente il ragionamento va fatto regola per regola, e per quelle comuni alla logica proposizionale il ragionamento è del tutto analogo a quello già visto, quindi ci concentreremo sulle regole di quantificazione. Iniziamo con \(\forall E\). Allora avremo un albero della forma:

INSERISCI ALBERO

Con \(t\) libero oer \(x\) in \(\psi\). Per ipotesi induttiva abbiamo che \(\Gamma \models \forall x, \psi\), e quindi per ogni \(M,\sigma\) si ha
\[M,\sigma\models \Gamma\implies M,\sigma\models \forall x, \psi\iff M,\sigma[x\mapsfrom d]\models\psi \]
per ogni \(d \in D\). Ovviamente \(\den{t}[M][\sigma]\in D\), e quindi \(M,\sigma\left[\mapsfrom \den{t}[M][\sigma]\right]\models \psi\). Ora dal teorema di sostituzione, se \(t\) è libero per \(x\) in \(\psi\), allora \(M,\sigma\models \subs{\psi}{x}{t}\). Essendo che \(t\) è libero per \(x\) in \(\psi\) per ipotesi, allora \(M,\sigma\models \phi\), ovvero \(\Gamma \models \phi\).